% Cover Letter - Healthcare Analytics (Elsevier)
\documentclass[12pt]{article}
\usepackage[margin=1in]{geometry}
\usepackage{times}

\begin{document}

\begin{flushleft}
[Date]\\[2em]
\end{flushleft}

The Editor-in-Chief\\
\textit{Healthcare Analytics}\\
Elsevier\\[2em]

Dear Editor-in-Chief,\\[1em]

We are pleased to submit our original research article entitled \textbf{``Continuum Memory Architecture (CMA) for Clinical Search: A Simulation-Based Evaluation of Intent-Aware Retrieval Using Complex Public Clinical Data''} for consideration for publication in \textit{Healthcare Analytics}.

Electronic health record (EHR) chart review is a non-linear, multi-turn task in which clinicians pivot abruptly between topics. Conventional session-based retrieval embeds prior queries uniformly, so stale context contaminates rankings after a topic shift and increases cognitive load. In this study we evaluate the Continuum Memory Architecture (CMA), which represents session intent as a trajectory on the Riemannian manifold of symmetric positive definite (SPD) matrices, suppresses stale context with a geodesic-shift-aware interference gate, and forecasts the next latent intent with a Joint Embedding Predictive Architecture (JEPA) to prefetch relevant documents.

In a randomized three-arm crossover simulation on clinical vignettes derived from publicly available EHR data (MIMIC-III/IV, eICU, HiRID), CMA reduced time-to-correct-information by 17.0\%, query latency by 20.8\%, and simulated cognitive load by 18.6\% while preserving perfect task completion, relative to TF-IDF and BM25 session-based baselines. These gains were stable and statistically significant across a scaling study from 1 to 10{,}000 vignettes (3--30{,}000 sessions), indicating that intent-aware session management improves the speed and efficiency of clinical chart review at scale.

This work fits the scope of \textit{Healthcare Analytics} by addressing an operational problem---efficient clinical information retrieval during chart review---with predictive analytics and machine learning, and by providing a reproducible simulation-based methodology for evaluating retrieval systems in healthcare organizations. It is an original contribution that has not been published previously and is not under consideration elsewhere. All authors have read and approved the submitted version and agree to its submission. The authors declare no competing interests.

We thank you for your consideration and look forward to your response.

\vspace{2em}
Sincerely,\\[2em]

Hemanth Manchabale Papachappa (corresponding author)\\[1em]

Aniruddha Ganguly

\end{document}
